\section{The calibrated predictor and its closed forms}\label{sec:method}

This section states the method under study, CrusherBayes as released in ChancaDEM~\productrelease{} and reported
in~\cite{cb002}, in the form the theory needs: the full Bayesian model, its predictive given the scalars, a linearization in
which the scalars are learned in closed form, and the three discrepancy priors compared in the study. Nothing in the
calibrated model is changed; the closed forms are exact consequences of it, and the full sampler is run beside them
(Section~\ref{sec:results-sampler}) to measure what they leave out.

\subsection{The model}

For survey $i$ and output $t\in\{\mathrm{tph},\mathrm{kW}\}$ (throughput in t/h, power in kW),
\begin{equation}
y_{t,i}=f_t(\boldsymbol\theta;\mathbf x_i)+\tilde{\mathbf z}_i^{\top}\mathbf d_t+\varepsilon_{t,i},\qquad
\mathbf d_t\sim N(\mathbf 0,\tau_t^2\mathbf I),\quad \varepsilon_{t,i}\sim N(0,\sigma_t^2),
\label{eq:model}
\end{equation}
where $\mathbf x_i=(c_i,f_{80,i})$ is the operating point (the closed-side setting and the feed size, both in mm), $f_t$ is the
engine, $\boldsymbol\theta$ are its eight fitted scalars, and the discrepancy regressors are
\begin{equation}
\tilde{\mathbf z}_i=\Big(1,\ \frac{c_i-\bar c}{\varsigma_c},\ \frac{f_{80,i}-\bar f}{\varsigma_f}\Big),\qquad \mathbf Z=[\tilde{\mathbf z}_1,\dots,\tilde{\mathbf z}_n]^{\top},
\label{eq:design}
\end{equation}
standardized with the mean and the population standard deviation of the training surveys. The engine is a Whiten
classification and breakage model~\cite{whiten,napier} of an HP500 secondary cone crusher with the calibration of Rocha et
al.~\cite{rocha}: the capacity $Q=Q_{\mathrm{ref}}(c+17)/38$ at the machine's fixed speed and throw, the Bond draw $P_d=W Q$
with the specific energy $W=10\,W_i(P_{80}^{-1/2}-F_{80}^{-1/2})$ in kWh/t~\cite{bond}, where $F_{80}$ is the 80\pct{}
passing size of the feed the engine reconstructs from the survey's $f_{80}$ (a Rosin-Rammler law read on the engine's sieve
grid, so $84.86$~mm for the 88~mm of survey~1), and the current-based power
$P=\zeta P_d+P_n$. The eight scalars are the capacity base $Q_{\mathrm{ref}}$, the work index $W_i$, the power slope $\zeta$,
the no-load power $P_n$, two multipliers $s_{K_1},s_{K_2}$ of the classification sizes, the classification exponent $K_3$ and a
shift $\Delta t_{10}$ of the breakage index. Their priors are truncated normals on the natural logarithm of the positive
scalars and on the value of the others, centred on the engine's shipped values (Table~\ref{tab:scalars}); the noise scales
have $\ln\sigma_t\sim N(\ln\sigma_{t,0},1)$ with $\sigma_{\mathrm{tph},0}=60$~t/h and $\sigma_{\mathrm{kW},0}=15$~kW; and the
discrepancy scales are $\tau_{\mathrm{tph}}=140$~t/h and $\tau_{\mathrm{kW}}=36$~kW, about 12\pct{} of the typical level of
each output. Given $\boldsymbol\theta$ and $\sigma_t$, the discrepancy is integrated out exactly,
\begin{equation}
\mathbf y_t\mid\boldsymbol\theta,\sigma_t\sim N\big(\mathbf f_t(\boldsymbol\theta),\ \sigma_t^2\mathbf I+\tau_t^2\mathbf Z\mathbf Z^{\top}\big),
\label{eq:marginal}
\end{equation}
and the ten-dimensional posterior of the eight scalars and the two noise scales is sampled by a differential-evolution
ensemble sampler~\cite{emcee,tb06,tbv08}. One posterior evaluates the engine about $1.26\times10^6$ times on nine surveys; at
\engineMs~ms per evaluation it costs about \cbCpuHours~CPU-hours.

\begin{table}[!t]
\caption{The eight fitted scalars of the engine on the HP500 lane, with their shipped value (the prior centre), the prior
standard deviation on the sampled scale, and the posterior mean of the fit to all ten surveys, which is the truth of the
virtual plants of Section~\ref{sec:design}.}
\label{tab:scalars}
\centering\footnotesize
\resizebox{\columnwidth}{!}{\begin{tabular}{@{}lrrr@{}}
\toprule
Scalar & Shipped & Prior sd & Plant truth \\
\midrule
$Q_{\mathrm{ref}}$ (t/h) & 647.00 & log 0.25 & 646.346 \\
$W_i$ (kWh/t) & 9.00 & log 0.17 & 8.387 \\
$\zeta$ & 1.30 & log 0.23 & 1.144 \\
$P_n$ (kW) & 110.00 & 50.00 & 86.353 \\
$s_{K_1}$ & 1.00 & log 0.15 & 1.068 \\
$s_{K_2}$ & 1.00 & log 0.15 & 1.051 \\
$K_3$ & 2.30 & 0.50 & 2.227 \\
$\Delta t_{10}$ (points) & 0.00 & 4.00 & -0.106 \\
\bottomrule
\end{tabular}
}
\end{table}

\subsection{The predictive given the scalars}

Given the scalars and the noise scale, the posterior predictive of a held-out survey at $\mathbf x_h$ is Gaussian. By the
push-through identity $\mathbf Z^{\top}(\sigma^2\mathbf I+\tau^2\mathbf Z\mathbf Z^{\top})^{-1}=(\sigma^2\mathbf I+\tau^2\mathbf
G)^{-1}\mathbf Z^{\top}$, with $\mathbf G=\mathbf Z^{\top}\mathbf Z$ and $\lambda_t=\sigma_t^2/\tau_t^2$, its mean and variance are
\begin{align}
\hat y_{t}(\mathbf x_h)&=f_t(\boldsymbol\theta;\mathbf x_h)+\tilde{\mathbf z}_h^{\top}(\mathbf G+\lambda_t\mathbf I)^{-1}\mathbf Z^{\top}\big(\mathbf y_t-\mathbf f_t(\boldsymbol\theta)\big),\label{eq:cbmean}\\
s_t^2(\mathbf x_h)&=\sigma_t^2\Big(1+\tilde{\mathbf z}_h^{\top}(\mathbf G+\lambda_t\mathbf I)^{-1}\tilde{\mathbf z}_h\Big).\label{eq:cbvar}
\end{align}
The mean~\eqref{eq:cbmean} is the engine corrected by a ridge regression of its residuals on the discrepancy regressors; the
posterior covariance of the discrepancy coefficients is $\sigma_t^2(\mathbf G+\lambda_t\mathbf I)^{-1}$, which
gives~\eqref{eq:cbvar}. Proposition~3 of~\cite{cb002} rewrites~\eqref{eq:cbmean} as a ridge regression of the survey on
$\mathbf Z$ toward the engine's best plane, plus the engine's departure from that plane at $\mathbf x_h$; Section~\ref{sec:theory}
computes its risk. The full method mixes the Gaussian predictives of~\eqref{eq:cbmean} and~\eqref{eq:cbvar} over the
posterior of $(\boldsymbol\theta,\sigma_t)$: its mean is the average of~\eqref{eq:cbmean}, and its band is a quantile of the
mixture, whose variance adds the spread of~\eqref{eq:cbmean} over the posterior to the average of~\eqref{eq:cbvar}. The closed
form fixes the scalars and the noise. Two fixings are used. The oracle fixing sets the scalars to the
truth of the virtual plant, or to the scalars that bring the engine closest to the plant in the $L_2$ sense of Tuo and
Wu~\cite{tuowu} when the plant is not an engine, and the noise to the plant's; it is the best the calibrated predictor can do
when the scalars are perfectly learned. The nominal fixing sets the scalars to their shipped values, the engine as it leaves
the factory with the discrepancy fitted around it.

\subsection{Learning the scalars in a linearized model}\label{sec:method-linear}

Between the fixed scalars and the full posterior lies a model in which the scalars are learned exactly. Linearize the engine
at a point $\boldsymbol\theta_c$ on the sampled scale, $\mathbf f_t(\boldsymbol\theta)\approx\mathbf f_{t,c}+\mathbf J_t\mathbf u$ with
$\mathbf u=\boldsymbol\theta-\boldsymbol\theta_c$ and $\mathbf J_t$ the Jacobian on the training surveys (central differences
with step $10^{-3}$ on the sampled scale). Then~\eqref{eq:model} becomes a Gaussian linear model in $\mathbf w=(\mathbf u,\mathbf
d_t)$ with features $\mathbf W_t=[\mathbf J_t,\mathbf Z]$,
\begin{equation}
\mathbf y_t-\mathbf f_{t,c}=\mathbf W_t\mathbf w+\boldsymbol\varepsilon_t,\qquad \mathbf w\sim N\big(\mathbf 0,\ \mathrm{blockdiag}(\boldsymbol\Lambda^{-1},\mathbf K_d^{-1})\big),
\label{eq:linear}
\end{equation}
with $\boldsymbol\Lambda$ the diagonal prior precision of the scalars from Table~\ref{tab:scalars} (the truncation dropped) and
$\mathbf K_d=\tau_t^{-2}\mathbf I$ for the prior of~\eqref{eq:model}. The posterior of $\mathbf w$ is Gaussian with precision
$\mathbf P_t=\mathbf W_t^{\top}\mathbf W_t/\sigma_t^2+\mathrm{blockdiag}(\boldsymbol\Lambda,\mathbf K_d)$, and with the held-out row
$\mathbf w_h=(\mathbf j_h,\tilde{\mathbf z}_h)$ the predictive is
\begin{equation}
\hat y_t(\mathbf x_h)=f_{t,c}(\mathbf x_h)+\mathbf w_h^{\top}\mathbf P_t^{-1}\mathbf W_t^{\top}(\mathbf y_t-\mathbf f_{t,c})/\sigma_t^2,\qquad
s_t^2(\mathbf x_h)=\sigma_t^2+\mathbf w_h^{\top}\mathbf P_t^{-1}\mathbf w_h .
\label{eq:linpred}
\end{equation}
The two outputs share $\mathbf u$; each output is fitted separately in~\eqref{eq:linpred}, which keeps the comparison per
output and leaves the information that one output carries about the other's scalars to the full sampler and to the design
study, where the joint form~\eqref{eq:eig} is used. The linearization point is the prior centre, the shipped engine, so the
learned scalars start where the full method starts. It is a first-order expansion of the engine at the prior centre, not of
the posterior at its mode; with the noise scales fixed, the truncation of the priors dropped and each output fitted alone,
\eqref{eq:linpred} is the exact posterior predictive of that per-output linear model, and it differs from the released
posterior on all three counts even for an engine linear in its scalars. The earlier paper checked a linearization at the
posterior mean against the sampler on the ten surveys and found the contraction of thirteen of fourteen parameters
reproduced~\cite{cb002}.

\subsection{The predictive band of each form}

The three forms of the method carry different shares of the uncertainty, and their bands differ accordingly. With the scalars
and the noise fixed, the band $\hat y_t\pm1.2816\,s_t$ of~\eqref{eq:cbmean} and~\eqref{eq:cbvar} carries the noise and the
uncertainty of the three discrepancy coefficients. The linearized form~\eqref{eq:linpred} adds the uncertainty of the scalars
through $\mathbf w_h^{\top}\mathbf P_t^{-1}\mathbf w_h$, which includes their correlation with the discrepancy. The full method
mixes the conditional predictives over the posterior of the scalars and the noise scale, so its predictive variance is
\begin{equation}
\mathrm{Var}(y_{t,h}\mid\mathcal D)=\mathbb E\big[s_t^2(\mathbf x_h)\mid\mathcal D\big]+\mathrm{Var}\big[\hat y_t(\mathbf x_h)\mid\mathcal D\big],
\label{eq:mixvar}
\end{equation}
the average conditional variance plus the spread of the conditional means, and its band is a pair of quantiles of the mixture,
which is not Gaussian when the posterior of the noise scale is skewed. A band built from the closed forms is therefore narrower
than the full method's whenever the scalars or the noise scale are uncertain; the study measures how much narrower on the
verification cells. The bands of the empirical comparators are of a different kind: a Student-$t$ prediction interval is exact
for a linear model with Gaussian noise and says nothing about the plant's curvature, which it absorbs into the residual variance.
Corollary~\ref{cor:coverage} gives the exact coverage of the closed-form band under any plant, which is what lets the study
separate the effect of the bias on coverage from that of the width.

\subsection{Three discrepancy priors}\label{sec:method-priors}

The prior of~\eqref{eq:model} is the one of Kennedy and O'Hagan~\cite{ko} reduced to a plane: the discrepancy may take any
direction of the regressors, including the directions in which the engine's scalars move the output, and the scalars and
the discrepancy are told apart only by their priors. Two published alternatives restrict this freedom.

The orthogonal prior of Plumlee~\cite{plumlee} forbids the discrepancy the directions of the engine's sensitivities, so that
a scalar keeps the meaning of the value that makes the engine closest to the plant. For a Gaussian process discrepancy the
construction subtracts from the covariance its projection on the model's gradient (Plumlee and Joseph~\cite{plumleejoseph},
Theorem~1; Tuo~\cite{tuo}, Corollary~3.6). For the plane of~\eqref{eq:model} the construction is finite: with $\mathbf J^{+}$ the
Moore-Penrose inverse of the training Jacobian, the discrepancy regressors become
\begin{equation}
\mathbf Z_\perp=(\mathbf I-\mathbf J\mathbf J^{+})\mathbf Z,\qquad \tilde{\mathbf z}_{h,\perp}=\tilde{\mathbf z}_h-(\mathbf J^{+}\mathbf Z)^{\top}\mathbf j_h,
\label{eq:orth}
\end{equation}
orthogonal to every column of $\mathbf J$ over the training surveys, and the held-out row is mapped by the same linear map,
which is the empirical form of an orthogonality with respect to the distribution of the operating points.

The scaled prior of Gu and Wang~\cite{guwang} conditions the discrepancy on a small mean square over the input domain, so that
the calibrated engine alone predicts well. For the plane, the mean square over the training surveys is
$\frac1n\|\mathbf Z\mathbf d\|^2=\mathbf d^{\top}(\mathbf G/n)\mathbf d$, and the scaled prior has precision
\begin{equation}
\mathbf K_d=\tau^{-2}\big(\mathbf I+c\,\mathbf G/n\big),
\label{eq:scaled}
\end{equation}
with $c=1$ in the study. Since the regressors are standardized, $\mathbf G/n$ has a unit diagonal, so for uncorrelated
setting and feed size the scaled prior halves the prior variance of every coefficient; with a correlation $\varrho$ between
them it divides the prior variance by $2+\varrho$ along the direction in which the two regressors move together and by
$2-\varrho$ along the direction in which they move apart. The three priors are thus three choices of $\mathbf K$ and of the features in~\eqref{eq:linear}, and
Proposition~\ref{prop:priors} states what the choice changes.

\subsection{Cost and the evaluation algorithm}

The closed forms cost a few small matrix operations once the engine and its Jacobian are known at the operating points. The
study evaluates the engine and its Jacobian once on a pool of 4,000 operating points per window, so that a replicate costs
$O(n\,p^2+p^3)$ arithmetic with $p\le11$ features, against about $1.26\times10^6$ engine evaluations for one posterior of the
full method. Algorithm~\ref{alg:replicate} is one replicate of the study: it draws a design and a noise realization, computes
every predictor on the same data, and records the errors, the band hits and the theoretical risk of Theorem~\ref{thm:risk} for
that design, so that the simulation checks the theory replicate by replicate. The seeds are derived from the cell and the
replicate index, so a result does not depend on the order or the number of processes.

\begin{algorithm}[!t]
\caption{One replicate of a study cell: virtual plant $\mu_t$, noise $\sigma_t$, training window, test window, size $n$}
\label{alg:replicate}
\begin{algorithmic}[1]
\Require pools with the plant's mean, the engine and its Jacobian at $\boldsymbol\theta_c$ and at the oracle scalars; seed $(\text{cell},\text{replicate})$
\State draw $n$ training points and 200 test points from the pools, without replacement
\State $y_{t,i}\gets\mu_t(\mathbf x_i)+\sigma_t\xi_i$ with $\xi_i\sim N(0,1)$, for training and test points
\State standardize on the training points; build $\mathbf Z$, $\mathbf G$ and the test rows by~\eqref{eq:design}
\If{$\mathrm{cond}(\mathbf G)>10^{10}$} redraw the design and count it \EndIf
\State oracle and nominal fixings: mean and variance by~\eqref{eq:cbmean} and~\eqref{eq:cbvar}
\State learned scalars, three priors: features by~\eqref{eq:linear},~\eqref{eq:orth}; $\mathbf K_d$ by~\eqref{eq:scaled}; predictive by~\eqref{eq:linpred}
\State comparators: least squares on $\mathbf Z$ and on $(1,z_c)$, the engine nominal, oracle and refit, the constant
\State bands: Gaussian $\pm1.2816\,s$; Student-$t$ prediction intervals for the regressions
\State record squared and absolute percentage errors, band hits and widths per predictor
\State record the risk of Theorem~\ref{thm:risk} for this design, averaged over the test points
\State \Return the replicate's records
\end{algorithmic}
\end{algorithm}

The full sampler is run, at a reduced size of 20 walkers and 4,000 steps of which 1,000 are discarded, on two tasks the closed
forms cannot do: simulation-based calibration, which checks that the sampler returns the posterior of the model, and a
verification on three cells in which the full posterior replaces the closed forms on the same replicates. The reduced size
keeps more than 50 integrated autocorrelation times after the burn-in, the convergence criterion of~\cite{cb002}, on the
posterior of the ten surveys, whose largest autocorrelation time is \tauTen~steps.
